\section{Marco Teórico}

El presente estudio se enmarca en el campo de las ciencias naturales, específicamente en la física, dentro del análisis del decaimiento radiactivo. Este fenómeno describe la transformación espontánea de núcleos atómicos inestables a nivel microscópico a lo largo del tiempo, así como sus implicaciones en distintos ámbitos, incluyendo el impacto ambiental, económico y energético derivado de una gestión inadecuada de materiales radiactivos.

El marco teórico se fundamenta principalmente en la ecuación de decaimiento exponencial:

\[
N(t) = N_0 e^{-\lambda t}
\]

donde \(N(t)\) representa el número de núcleos restantes en el tiempo \(t\), \(N_0\) es la cantidad inicial de núcleos y \(\lambda\) es la constante de decaimiento característica de cada material.

Eventos históricos como el desastre en la central nuclear de Chernóbil en 1986 y el accidente en la central de Fukushima en 2011 han evidenciado las graves consecuencias a largo plazo asociadas a la liberación de material radiactivo en el medio ambiente.

Las radiaciones emitidas, tales como \(\alpha\), \(\beta\) y \(\gamma\), pueden permanecer en el suelo, el aire y las fuentes de agua, generando contaminación persistente. Entre los materiales involucrados se encuentra el \ce{^{238}U}, el cual posee una vida media del orden de $4.5\times10^9$ años, lo que evidencia la magnitud temporal de sus efectos.

\subsection{Tipos de radiación}

Como se mencionó anteriormente, existen diferentes emisiones mediante las cuales los núcleos inestables buscan alcanzar una configuración más estable. Los tres tipos principales son:

\begin{itemize}

    \item \textbf{Radiación Alfa ($\alpha$):} Consiste en la emisión de un núcleo de helio, formado por dos protones y dos neutrones. Posee una gran masa (más de 7000 veces la de una partícula beta) y carga eléctrica, lo que le da un alto poder de ionización pero un alcance muy corto en la materia, siendo incapaz de atravesar incluso la piel humana.
    
    \item \textbf{Radiación Beta ($\beta$):} Se produce por la transformación de un neutrón en protón (emisión $\beta^-$), liberando un electrón a velocidades relativas junto con un neutrino. Su capacidad para atravesar la materia es intermedia, mayor que las partículas alfa pero significativamente menor que los rayos gamma.
    
    \item \textbf{Radiación Gamma ($\gamma$):} Es una forma de radiación electromagnética compuesta por fotones de alta energía, ocurre cuando un núcleo pasa de un estado excitado a uno de menor energía sin cambiar su composición química. Este tipo de radiación es la más penetrante y difícil de detener, ya que requiere materiales densos como el plomo para su blindaje.
\end{itemize}

\subsection{Permanencia en el Medio Ambiente}

La persistencia de la radiación en el medio ambiente no es uniforme, ya que depende de diversos factores físicos y químicos. Entre ellos, la vida media de los radionúclidos es determinante, pues establece el tiempo característico durante el cual un material permanece radiactivamente activo.

En condiciones naturales, la radiactividad forma parte del entorno y contribuye, por ejemplo, al calor interno de la Tierra. No obstante, los accidentes nucleares introducen al ambiente radionúclidos específicos en concentraciones elevadas, alterando significativamente el equilibrio natural. Un ejemplo relevante es el isótopo \ce{^{131}I}, que, a pesar de tener una vida media relativamente corta (aproximadamente 8 días), representa un riesgo biológico considerable debido a su alta actividad inicial y su fácil incorporación en organismos vivos.

Asimismo, eventos como los ocurridos en Chernóbil y Fukushima han evidenciado que la liberación de material radiactivo puede generar zonas inhabitables durante largos periodos de tiempo, afectando tanto el área del accidente como las regiones circundantes. En este contexto, el estudio de los efectos a largo plazo de la radiación —incluyendo impactos ambientales, sanitarios, económicos y energéticos— resulta fundamental para evaluar la viabilidad y sostenibilidad de la energía nuclear como alternativa energética.
